\documentclass[11pt,a4paper]{article}
\usepackage[margin=23mm]{geometry}
\usepackage{amsmath,amssymb,amsthm,booktabs}
\usepackage[hidelinks]{hyperref}
\newtheorem{proposition}{Proposition}
\newcommand{\E}{\mathbb E}
\newcommand{\F}{\mathcal F}
\title{Economic authorization and statistical certification}
\author{PARK--REVIVE research design note}
\date{18 October 2026}
\begin{document}
\maketitle
\noindent\textbf{Status.} Proposed extension following the economic-selection
discussion. This note supplies a target distinction and an experiment design;
it reports no new backtest, simulation result, power guarantee or frozen
numerical protocol. The six-method, 500-market study already running retains
its original configuration, seed membership and analysis plan.

\section{The economic object is a learned policy}
For a common pre-trade allocation $b$, define
\[
 C(w,b)=c\|w-b\|_1+\frac{\eta}{2}(w-b)^\top\Lambda(w-b),
 \qquad \Lambda\succeq0,
\]
and realized quadratic utility
\[
 u(w,b;A)=w^\top A-\frac{\gamma}{2}(w^\top A)^2-C(w,b).
\]
With common true moments and the same feasible constraints, the optimized
value
\[
 V_t(S;b)=\sup_{\operatorname{supp}(w)\subseteq S}
             \E[u(w,b;A_{t+1})\mid\F_t]
\]
cannot decrease when $S$ expands: the larger feasible set includes the
smaller one. This observation assumes that adding a permitted feature does
not itself impose a separate fixed participation charge. It does not imply
monotonicity of the out-of-sample values of estimated policies:
\[
 \Delta_t(B,S)=
 \E\!\left[
 u(\widehat w_t(B),b^B_t;A_{t+1})
 -u(\widehat w_t(S),b^S_t;A_{t+1})\mid\F_t\right].
\]
Estimation error, changes in allocation and the two cost states matter.
Even estimated optimized training value is monotone under nested supports;
its positive increment alone is therefore weak protection against overfit.

The current full-study \texttt{Greedy} method is not a pure training-only
economic selector. It fits allocations using the training history, but
eligibility requires a positive accumulated OOS shadow-score mean, and
annual priority ranks eligible moves by that mean. A method that initiates
exposure from training estimates and estimation uncertainty would be an
additional comparator.

\section{Separate permission to trade from the tested null}
A training-based forecast can authorize provisional exposure subject to a
declared risk limit. An uncertainty adjustment must specify its estimation
model, horizon and calibration before its new evaluation. A working posterior,
bootstrap percentile or confidence bound does not acquire an OOS guarantee
merely by being used as an economic gate.

For a predictable pair of shadow policies, an e-process addresses the
sequential null that its score has nonpositive conditional mean throughout
the comparison. At individual level $a$, crossing $1/a$ controls rejection
under that null, assuming the score moment contract. The threshold is $20$
when $a=.05$ for one comparison; using it repeatedly across features and
episodes does not by itself control a global false-authorization probability.

Under the price-taking simulation contract, predictable provisional
investment need not invalidate a shadow e-process that retains exactly the
same policies, cost histories, filtration and valid moment bounds.
Changing the reference, candidate, weights or cost state changes the tested
object. The validity of the construction must then be checked for that
object; previous capital cannot simply be relabeled as evidence for a
different comparison. Neither a shadow rejection nor an economic forecast
guarantees a positive future deployment gain in a changing market.

\subsection*{Certification against cash does not certify an increase}
Suppose the current allocation is $qw$, $0<q<1$, and compare increasing to
$w$ against retaining $qw$ from that same state. Put
\[
 m_1=\E[w^\top A\mid\F_t],\quad
 m_2=\E[(w^\top A)^2\mid\F_t],\quad
 L_w=\|w\|_1,\quad K_w=w^\top\Lambda w.
\]
The conditional utility gain from the increase is
\begin{align}
 \Delta_{\mathrm{increase}}
 &=(1-q)m_1-\frac{\gamma}{2}(1-q^2)m_2
       -c(1-q)L_w-\frac{\eta}{2}(1-q)^2K_w. \label{eq:increase}
\end{align}
The gain from entering $w$ from cash is
\[
 \Delta_{\mathrm{entry}}
 =m_1-\frac{\gamma}{2}m_2-cL_w-\frac{\eta}{2}K_w.
\]
Consequently
\[
 \Delta_{\mathrm{increase}}
 =(1-q)\Delta_{\mathrm{entry}}
   -\frac{\gamma}{2}q(1-q)m_2
   +\frac{\eta}{2}q(1-q)K_w.
\]
For example, with no costs, $\gamma=1$, $m_1=.6$, $m_2=1$ and $q=.5$,
full exposure has expected utility $.1$, half exposure has $.175$, and
increasing exposure loses $.075$. These are admissible moments: a gross
return with mean $.6$ and variance $.64$ has second moment $1$.

Thus a certificate for full exposure versus cash does not establish that
increasing a provisional position is beneficial. A favorable economic
forecast for the increase must use \eqref{eq:increase}, or its general
state-dependent analogue. If the increase is to carry its own statistical
certification, its candidate/reference score needs an appropriate e-process
and error allocation. Keeping the original cash comparison is useful as
a diagnostic, with its original interpretation.

\section{A comparison that isolates the decision to wait}
The full-study Tail-minus-Greedy contrast compares complete adaptive funds.
It includes different selected supports, allocation states and later moves,
so it does not isolate the effect of waiting. Add a controlled mechanism
experiment with the following common objects:
\begin{enumerate}
 \item Use the same training data and training-only candidate-selection rule
 in both arms. Define the candidate allocation path $w_t$ and all subsequent
 fitting dates identically, without using either arm's wealth or trading
 state to change that common path.
 \item Define a training uncertainty rule and a provisional fraction
 $q\in[0,q_{\max}]$, with $q_{\max}<1$, before examining new results.
 A rejected economic proposal has $q=0$ and remains in the experiment.
 \item Run one common shadow candidate-versus-cash e-process using the
 declared score and its own shadow costs. Define a common, predictable
 full-exposure date from its crossing and any common economic gate.
 A decision made after a return is first tradable on the following date.
 \item The waiting arm holds cash before that date. The provisional arm
 holds $qw_t$ before it. Both target $w_t$ afterwards. Each pays its own
 actual costs, including the different transition costs at the increase.
 The original shadow certificate remains a candidate-versus-cash
 certificate in both arms; it is not labeled a certificate of $q\to1$.
 \item Retain the complete evaluation horizon, including noncrossings,
 rejected training proposals and losses. Use paired markets with common
 shocks, and report separately the provisional losses and later outcomes.
\end{enumerate}
This experiment identifies an exposure-timing contrast for the specified
common candidate path and evidence clock. It does not establish the
performance of a general adaptive feature-selection algorithm. A separate
extension could test the increase directly, but that changes the evidence
clock and needs its own design and guarantees.

\begin{proposition}[Realized utility accounting]
Let $\theta^P_t,\theta^W_t$ be the provisional and waiting exposure fractions,
$X_t=w_t^\top A_t$ the common full-target gross return, and $C^P_t,C^W_t$
their actual trading costs. Their realized utility difference over $H$ dates
is exactly
\[
 D_H=\sum_{t=1}^H
 \left[(\theta^P_t-\theta^W_t)X_t
 -\frac{\gamma}{2}\bigl((\theta^P_t)^2-(\theta^W_t)^2\bigr)X_t^2
 -(C^P_t-C^W_t)\right].
\]
\end{proposition}
\begin{proof}
Substitute $\theta^P_tw_t$ and $\theta^W_tw_t$ into the utility definition
and subtract, retaining each arm's own cost state.
\end{proof}
With $q$ fixed before the common deployment date $T$ and full exposure in
both arms afterwards, the gross-return and quadratic-risk terms cancel for
$t\ge T$. Under the repository's cost convention, whose state is the previous
target allocation, costs also coincide after the deployment date. Costs
before and at $T$ remain arm-specific. If $T>H$, including $T=\infty$, the
pre-deployment contrast extends through all $H$ dates. Wealth compounding is
a distinct outcome and is not this additive utility identity.

\section{Outputs and implementation boundary}
The primary mechanism outcome should be the paired net utility difference,
with expected-benefit, realized gross-return, quadratic-risk and trading-cost
components retained separately. Risk outcomes should include negative
provisional utility, downside outcomes, drawdowns and insolvency under the
same wealth convention; none is a type-I error guarantee. Report waiting
durations and noncrossings without conditioning performance on certification.

Before implementation and new evaluation, fix the uncertainty estimator and
calibration, candidate-selection rule, provisional cap, investment horizon,
review schedule, complete seed set, error allocation, outcome definitions
and Monte Carlo analysis. The design follows observed pilot results and the
user's proposal; it is not a pre-pilot registration. Reuse of existing market
paths would be a paired extension, not additional independent market evidence.

No new numerical superiority claim follows from this note. In particular,
the question \emph{what benefit is forgone by waiting, and what downside is
avoided?} requires the new controlled comparison; the existing adaptive fund
contrast supplies complementary total-policy performance.

\subsection*{Implementation references inspected}
\noindent\path{research/park_revive_full/selection.py}:\\
\texttt{Selector.refresh}, \texttt{eligible}, \texttt{choose};

\noindent\path{research/park_revive_full/engine.py}:\\
OOS score updates and annual decisions.

\noindent\path{src/park_revive/policies.py}: previous-target cost state.

\noindent\path{src/park_revive/costs.py}: common utility and cost definitions.
\end{document}
